\documentclass[11pt,letterpaper]{article}
\usepackage[T1]{fontenc}
\usepackage[utf8]{inputenc}
\usepackage{lmodern,amsmath,amssymb,amsthm,mathtools}
\usepackage[margin=1in,headheight=14pt]{geometry}
\usepackage{microtype,booktabs,tabularx,array,enumitem,xcolor,fancyhdr}
\usepackage[hyphens]{url}
\usepackage[colorlinks=true,linkcolor=blue!45!black,citecolor=blue!45!black,urlcolor=blue!45!black]{hyperref}
\hypersetup{pdftitle={Faltings' theorem and intuitionistic primitive recursive arithmetic},pdfauthor={Research report prepared for Vladimir Reshetnikov}}
\setlist{topsep=4pt,itemsep=3pt}
\setlength{\emergencystretch}{2em}
\setlength{\parskip}{4pt}
\newcommand{\iPRA}{\mathsf{iPRA}}
\newcommand{\N}{\mathbb N}
\newcommand{\Z}{\mathbb Z}
\newcommand{\Q}{\mathbb Q}
\newcommand{\Good}{\operatorname{Good}}
\newcommand{\Pt}{\operatorname{Pt}}
\newcommand{\HB}{\operatorname{HB}}
\newcommand{\Comp}{\operatorname{Complete}}
\newtheorem{theorem}{Theorem}[section]
\newtheorem{lemma}[theorem]{Lemma}
\newtheorem{proposition}[theorem]{Proposition}
\newtheorem{corollary}[theorem]{Corollary}
\theoremstyle{definition}
\newtheorem{definition}[theorem]{Definition}
\newtheorem{remark}[theorem]{Remark}
\pagestyle{fancy}\fancyhf{}
\fancyhead[L]{\small Faltings' theorem and intuitionistic PRA}
\fancyhead[R]{\small Research and formalization audit}
\fancyfoot[C]{\thepage}
\begin{document}
\begin{titlepage}
\centering\vspace*{1.2cm}
{\LARGE\bfseries Faltings' theorem and intuitionistic\\[5pt]
primitive recursive arithmetic\par}
\vspace{0.6cm}
{\large A precise effectivity obstruction, constructive special cases,\\
and a checked Lean foundation\par}
\vspace{0.9cm}
{\large Prepared for Vladimir Reshetnikov\par}
\vspace{0.2cm}{November 15, 2026\par}
\vspace{0.8cm}
\begin{minipage}{0.91\textwidth}
\textbf{Status.} This report does not prove general Faltings in $\iPRA$.
It proves that a uniform $\iPRA$ proof of a positive finite-list or
height-bound formulation would yield a primitive recursive version of
effective Mordell, which the
sources checked here do not establish. This is not an independence proof,
and it does not rule out a weaker cardinality formulation.

The positive results are complete human-readable proofs of the finite-search
reductions, the four-point classification of the Fermat quartic, and
nonexistence for a simple hyperelliptic family. Three standalone Lean modules
verify the elementary search, logical, and arithmetic layers.
The general height statement is explicitly defined and left unproved.
No general finiteness axiom or proof placeholder is introduced.
\end{minipage}
\vfill
\begin{minipage}{0.91\textwidth}\small
\textbf{Repository snapshot.}
\href{https://github.com/VladimirReshetnikov/ProveIt/commit/d72c8b24637e342816dc92d11bef337748683e51}
{ProveIt, commit \texttt{d72c8b24637e342816dc92d11bef337748683e51}}.

\textbf{Verification.} Lean 4.32.0, standard library only; exact integer
and rational Python checks. Native Lean theorems are distinguished from
derivations in the repository's embedded $\iPRA$ calculus.
\end{minipage}
\end{titlepage}

\begin{abstract}
The requested general proof cannot be supplied by the arguments established
here. Arbitrary-matrix intuitionistic witnessing turns a uniform $\iPRA$
proof of positive rational-point finiteness into a primitive recursive
height bound. We make this implication precise for squarefree monic
quintics, give all arithmetic predicates and an explicit enumeration, and
prove the equivalence of finite-list and height witnesses. Cardinality
bounds and double-negated finiteness are separated from this conclusion.
As an unconditional substantial instance, explicit finite descent proves
that the Fermat quartic has exactly four rational points. The accompanying
Lean files compile without holes or custom axioms; the audited results
have no dependence on \texttt{Classical.choice}. The report specifies
the remaining object-theory work and the unresolved general mathematical
step rather than placing that step inside an assumed lemma.
\end{abstract}
\setcounter{tocdepth}{1}
\tableofcontents
\newpage

\section{The request and the actual research boundary}
Faltings' theorem says that if $K$ is a number field and $C/K$ is a
smooth, projective, geometrically integral curve with genus $g(C)\geq2$,
then $C(K)$ is finite. A number field is a finite extension of $\Q$;
the genus of a smooth projective curve is the dimension of its space of
regular differentials. These geometric conditions are part of the theorem.

Stoll's ICM 2026 account states that the known proofs are ineffective:
they do not give a general algorithm to determine the finite rational-point
set \cite{Stoll}. There are effective results under additional hypotheses.
Garcia-Fritz and Pasten obtain height bounds when a curve has enough
automorphisms relative to its Jacobian rank \cite{GFP}.
Alp\"oge and Lawrence give a general machine whose termination follows
from the Hodge, Tate, and Fontaine--Mazur conjectures \cite{AL}.
Neither supplies the unconditional weak-arithmetic proof requested here.

Recent explicit \emph{cardinality} bounds also matter. Yu, Yuan, and
Zhou give numerical bounds on the number of points, including
discriminant-dependent bounds for monic hyperelliptic equations
\cite[Theorem 1.7]{YYZ}. Such bounds do not themselves supply heights
or a stopping rule for enumerating the points.

\begin{center}
\begin{tabularx}{\textwidth}{@{}p{0.34\textwidth}X@{}}
\toprule
Claim & Status here\\ \midrule
General positive Faltings in $\iPRA$ & Not proved; the exact missing target is displayed below.\\
Unprovability or independence & Not claimed.\\
Bounded search and list/height equivalence & Fully proved; elementary Lean versions compile.\\
Positive $\iPRA$ proof gives a primitive recursive bound & A proved consequence of a cited external witnessing theorem.\\
Fermat quartic over $\Q$ & Complete constructive human argument; embedded Lean derivation remains to be constructed.\\
Empty even-degree hyperelliptic family & Complete sign argument; its arithmetic core is checked in Lean.\\
\bottomrule
\end{tabularx}
\end{center}

The original request admits more than one constructive reading.
A cardinality-bound statement is weaker than a positive finite-list
statement. The height-witness obstruction applies to the latter.
This report neither proves nor refutes general $\iPRA$ provability of
the former. Choosing a weaker formulation must be made explicit.

\section{Object theory and finiteness specifications}
\label{sec:setting}
We use standard first-order intuitionistic primitive recursive arithmetic.
Its language has $0$, successor, and symbols for specified primitive
recursive functions. Zero, successor, and projections are initial.
Composition and primitive recursion have their defining equations
\[
 h(\vec x)=f(g_1(\vec x),\ldots,g_m(\vec x)),\qquad
 r(0,\vec x)=a(\vec x),\quad
 r(n+1,\vec x)=b(n,r(n,\vec x),\vec x).
\]
There are successor injectivity and nonzero-successor axioms, intuitionistic
first-order logic with equality, and induction only for quantifier-free
formulas with parameters:
\[
 [\varphi(0,\vec a)\wedge
 \forall n(\varphi(n,\vec a)\to\varphi(n+1,\vec a))]
 \to\forall n\,\varphi(n,\vec a).                         \tag{I}
\]
This is the purpose of \texttt{PRAAx} and \texttt{PRA.qfInd} in the
inspected repository \cite{Repo}.

Finite loops, bounded search, integer arithmetic, gcd, and coded lists
are primitive recursive. Bounded quantifiers over primitive recursive predicates
are implemented by Boolean scans with proved invariants. They do not
authorize unrestricted induction. A parameter can remain free during
quantifier-free induction and be universally generalized afterwards.

For a primitive recursive point predicate $\Pt(c,p)$ and an integer
height $H$, define positive height finiteness by
\[
 \exists B\,\HB(c,B),\qquad
 \HB(c,B):=\forall p\,(\Pt(c,p)\to H(p)\leq B),              \tag{H}
\]
and positive list finiteness by
\[
 \exists L\,\Comp(c,L),\qquad
 \Comp(c,L):=\forall p\,(p\in L\leftrightarrow\Pt(c,p)).     \tag{L}
\]
Lists and membership are finitely coded. A cardinality bound instead says
\[
 \exists N\,\forall L\,
 [\operatorname{Nodup}(L)\wedge\operatorname{AllPoint}(c,L)
   \to |L|\leq N].                                      \tag{C}
\]
The predicates on $L$ are finite checks. The numeral $N$ need not reveal
the coordinates or heights of any point. Section \ref{sec:logic} explains
why (C) cannot generally be converted constructively into (H) or (L).

Quantification over all infinite sequences is not first-order $\iPRA$.
An assertion excluding injective infinite point sequences belongs either
to an ambient higher-order setting or to a separately specified arithmetic
schema. Neither is silently identified with (H).

\section{An explicit genus-two arithmetic target}
\label{sec:quintic}
Write
\[
 f_{\mathbf a}(X)=X^5+a_4X^4+a_3X^3+a_2X^2+a_1X+a_0,
 \quad \mathbf a\in\Z^5,\qquad
 \Good(\mathbf a):\Longleftrightarrow
 \operatorname{Res}(f_{\mathbf a},f'_{\mathbf a})\ne0.      \tag{1}
\]
The resultant is a fixed $9$ by $9$ Sylvester determinant: four shifted
rows of $(1,a_4,a_3,a_2,a_1,a_0)$ and five shifted rows of
$(5,4a_4,3a_3,2a_2,a_1)$, padded to length nine. It is an integer
polynomial in the five coefficients. Hence $\Good$ is decidable
and primitive recursive.

This test is squarefreeness over $\Q$. The matrix represents
$(U,V)\mapsto Uf+Vf'$ with $\deg U<4$, $\deg V<5$.
If $f,f'$ are coprime, a kernel vector satisfies $f\mid V$, so the
degree restriction gives $V=0$, then $U=0$.
If they have a nonconstant gcd $h$, the vector $(f'/h,-f/h)$ is a
nonzero kernel vector of the required degrees. Finite-dimensional
determinant elimination proves the claim. Polynomial Euclidean division
has decreasing degree, and rational operations in this fixed dimension
admit finite arithmetic coding.

\subsection{Canonical coordinates and the point predicate}
Use reduced fractions
\[
 x=u/v,\quad y=r/s,\quad v,s>0,\quad
 \gcd(|u|,v)=\gcd(|r|,s)=1.
\]
Define
\[
 F_{\mathbf a}(u,v)=u^5+a_4u^4v+a_3u^3v^2+
                    a_2u^2v^3+a_1uv^4+a_0v^5.
\]
The predicate $\Pt(\mathbf a;u,v,r,s)$ consists of those positivity
and gcd conditions and
\[
                 r^2v^5=s^2F_{\mathbf a}(u,v).            \tag{2}
\]
As $v,s$ are positive, clearing or restoring denominators proves that
(2) is exactly $y^2=f_{\mathbf a}(x)$. Set
\[
                 H(u,v,r,s)=\max(|u|,v,|r|,s).           \tag{3}
\]
All tests are primitive recursive; no equality of arbitrary real or
complex numbers is being decided.

\subsection{The geometric interpretation}
For $\Good(\mathbf a)$, glue the charts
\[
 y^2=f_{\mathbf a}(x),\qquad
 z^2=t(1+a_4t+a_3t^2+a_2t^3+a_1t^4+a_0t^5)              \tag{4}
\]
by $t=1/x$, $z=y/x^3$ where $x\ne0$.
On the first chart a singular point would satisfy $y=0$ and $f=f'=0$.
On the second chart at $t=0$, the unique point is $(0,0)$ and the
$t$ derivative of its defining equation is $-1$. Thus there is exactly
one smooth rational point at infinity.
The two monic quadratic charts give a finite morphism to $\mathbf P^1$,
so the resulting curve is projective. Over an algebraic closure,
$f_{\mathbf a}(x)$ has odd pole order five at infinity and cannot be
a square in the rational function field. Hence its quadratic function
field, and the smooth curve just constructed, are geometrically integral.

The double cover of $\mathbf P^1$ is branched at five simple finite
roots and at infinity, and has genus two. An explicit differential
calculation gives the same answer. At infinity the orders of $x,y,dx$
are $-2,-5,-3$. Anti-invariant regular differentials have the form
$r(x)\,dx/y$. A finite pole of $r$ would remain a pole, even at a
branch point where $dx/y$ is a unit differential; thus $r$ is a
polynomial. The orders at infinity allow exactly degrees zero and one.
The basis is $dx/y,x\,dx/y$. Every differential splits into invariant
and anti-invariant parts by $(\eta\pm\iota^*\eta)/2$, where
$\iota(x,y)=(x,-y)$. An invariant differential has the form $A(x)dx$.
At a branch point, a base pole of differential order $-m$, $m\geq1$,
pulls back with order $1-2m<0$; at an unbranched point its order is
unchanged. Thus a regular invariant differential descends regularly
to $\mathbf P^1$ and is zero.
This proves that the space has dimension two.

This is a semantic identification in ordinary algebraic geometry, not
a claim that the supplied Lean files develop a theory of schemes or
regular differentials. It agrees with the standard genus-two presentation
in \cite{Stoll}. The arithmetic point problem is precisely (2), with
one additional infinity code. In particular, projective nonemptiness
is automatic in this family and is not the claimed computational obstacle.

\subsection{The unproved positive statement}
A positive uniform theorem already implies the following subcase:
\[
 \boxed{\ 
 \forall\mathbf a\,
 \bigl(\Good(\mathbf a)\to
 \exists B\,\forall u,v,r,s\,
 [\Pt(\mathbf a;u,v,r,s)\to H(u,v,r,s)\leq B]\bigr).
 \ }                                                     \tag{F$_5$}
 \label{eq:target}
\]
Signed variables abbreviate natural-number codes.
In \texttt{Curves.lean} this proposition is defined as
\texttt{QuinticHeightGoal}. It is not declared as a theorem or axiom.

\section{Constructive enumeration and the equivalence of witnesses}
\label{sec:enum}
\begin{theorem}[Exact bounded enumeration]\label{thm:enum}
For every $\mathbf a$ and $B\in\N$ a primitive recursive construction
produces $E(\mathbf a,B)$ such that
\[
 p\in E(\mathbf a,B)\ \longleftrightarrow\
                 \Pt(\mathbf a,p)\wedge H(p)\leq B.       \tag{5}
\]
Its correctness admits quantifier-free induction, without assuming
squarefreeness.
\end{theorem}
\begin{proof}
List all tuples $-B\leq u,r\leq B$, $1\leq v,s\leq B$ in a fixed
lexicographic order and retain those passing (2) and the gcd tests.
There are exactly
\[
                         M(B)=(2B+1)^2B^2                \tag{6}
\]
raw candidates, none repeated. At $B=0$ the box is empty.

Let $E_j$ be the retained list after $j$ candidates. For a fixed tuple
$p$, let $\operatorname{Seen}(p,j)$ be its Boolean membership in that
prefix. The induction invariant is
\[
 \operatorname{Mem}(p,E_j)=
       [\operatorname{Seen}(p,j)\mathbin{\&}\Pt(\mathbf a,p)].
                                                               \tag{7}
\]
At zero both sides are false. If the next tuple differs from $p$,
neither side changes; if it equals $p$, both sides become the point
test. All cases are decidable and finite-list membership is a Boolean
scan. Thus the invariant is quantifier-free, with $p$ a parameter.
At $j=M(B)$, being seen means exactly membership in the component box.
Under the point predicate's positive-denominator conditions, this is
equivalent to the height condition (3). Generalize over $p$ after induction.
\end{proof}

\begin{theorem}[Positive lists and positive bounds]\label{thm:equiv}
For each $\mathbf a$, intuitionistically,
\[
 \exists L\,\Comp(\mathbf a,L)
             \quad\longleftrightarrow\quad
 \exists B\,\HB(\mathbf a,B).                              \tag{8}
\]
The witness transformations are primitive recursive.
\end{theorem}
\begin{proof}
Given $B$ and a proof of the global bound, take $L=E(\mathbf a,B)$.
Equation (5) gives soundness; the bound and (5) give completeness.

Conversely compute
$B_L=\max(\{0\}\cup\{H(p):p\in L\})$.
For a fixed $p$, the invariant
$\operatorname{Mem}(p,L_j)\to H(p)\leq B_{L_j}$ is quantifier-free
through the maximum scan. A new entry either is $p$, in which case the
maximum is at least its height, or preserves its old bound.
If $L$ is complete, every point occurs and is bounded by $B_L$.
For the empty list, completeness proves absence, so $0$ is a bound.
This proof never tries to decide an unbounded completeness assertion.
\end{proof}

\subsection{Numerical coding and an explicit finite-ball cap}
Encode integers by $\zeta(z)=2z$ for $z\geq0$ and
$\zeta(z)=-2z-1$ for $z<0$. Use the injective pairing
\[
                     \langle a,b\rangle=(a+b)^2+a.
\]
For $t=a+b$, its value lies between $t^2$ and $t^2+t<(t+1)^2$.
Bounded square root therefore recovers $t$, followed by
$a=n-t^2$, $b=t-a$. Codes with $n-t^2>t$ are malformed.
This is a primitive recursive decoder with a decidable validity test.

For the point code
$p=\langle\langle\zeta(u),v\rangle,\langle\zeta(r),s\rangle\rangle$,
height at most $B$ implies both inner pairs are at most
$A(B)=9B^2+2B$, and hence
\[
                        p\leq4A(B)^2+A(B).               \tag{9}
\]
Thus the geometric height balls have an explicit primitive recursive
code bound. It would not suffice to posit an arbitrary function named
height and assume that its balls could be enumerated.

Lists can use $\operatorname{nil}=0$ and
$\operatorname{cons}(a,L)=1+\langle a,L\rangle$.
A valid tail has smaller code, so fuel equal to the original code
suffices for primitive recursive decoding, length, membership, and
maximum. Malformed codes are totalized. These definitions complete the
arithmetic encoding needed in the preceding theorems.
In the numerical translation, malformed point codes make $\Pt$ false,
and malformed coefficient-tuple codes fail $\Good$. A decoder's default
zero value alone is never treated as a valid encoding. All integer
coefficient tuples in the typed presentation still have valid codes.

\section{A positive uniform proof would yield primitive recursive effective Mordell}
\label{sec:witness}
The residual matrix in \eqref{eq:target} is universally quantified.
A witnessing result restricted to decidable existential matrices would
not justify the next argument.

\begin{theorem}[External arbitrary-matrix witnessing]\label{thm:external}
If intuitionistic $\mathsf I\Sigma_1$ proves
$\forall x\exists y\,A(x,y)$ for an arbitrary arithmetic formula $A$,
there is a primitive recursive function $b$ such that
$\forall x\,A(x,b(x))$ is provable in its extension by primitive
recursive definitions. This stronger theory contains standard $\iPRA$.
\end{theorem}
This is Wehmeier's witnessing result, explicitly stated and reproved
by Ferreira and Marques \cite[p.~14 of the author manuscript]{FM}.
It is used as an external metatheorem, not as a Lean result of this
package. Its primitive-recursive-language presentation still has
$\Sigma_1$ induction; it is not identical to quantifier-free $\iPRA$.
We equivalently use the witnessing theorem in that definitional expansion:
the finitely many primitive recursive symbols in a given proof can be
introduced or eliminated by their verified definitions.

\begin{theorem}[Necessary computational consequence]\label{thm:necessary}
A single uniform $\iPRA$ proof of \eqref{eq:target} would yield a
primitive recursive $b:\Z^5\to\N$ satisfying, in the standard model,
\[
 \Good(\mathbf a)\to
 \forall p\,(\Pt(\mathbf a,p)\to H(p)\leq b(\mathbf a)).     \tag{10}
\]
It would yield a primitive recursive complete rational-point enumeration
for all squarefree monic quintics in this family.
\end{theorem}
\begin{proof}
The resultant test is decidable in $\iPRA$.
For fixed $\mathbf a$, if it succeeds, the hypothetical theorem gives
a witness $B$ and its bound proof. If it fails, choose $B=0$ and prove
the conditional bound from the contradictory goodness assumption.
Thus $\iPRA$ proves
\[
 \forall\mathbf a\,\exists B\,
                    [\Good(\mathbf a)\to\HB(\mathbf a,B)]. \tag{11}
\]
This is intuitionistic case analysis on a finite numerical test;
there is no case split on unbounded rational-point existence.

Code the five signed inputs by one natural number and include this
derivation in the stronger theory of Theorem \ref{thm:external}.
Apply its witnessing theorem to the arbitrary matrix
$\Good(\mathbf a)\to\HB(\mathbf a,B)$. The resulting primitive
recursive witness satisfies (10). Its correctness is obtained in that
stronger theory; we do not claim that this metatheorem already constructs
an $\iPRA$ proof of extracted correctness.

Finally compute $E(\mathbf a,b(\mathbf a))$ and append the unique point
at infinity. Theorems \ref{thm:enum} and \ref{thm:equiv} prove
completeness. The bounded loops and composition are primitive recursive.
\end{proof}

This is a necessary consequence, not a contradiction or an independence
result. A future effective proof is not excluded. It also distinguishes
uniformity from separate statements about fixed curves. Existence of a
maximum for each finite set does not give an algorithm for finding it,
nor does it prove that $\iPRA$ establishes each proposed maximum's
global correctness.

\section{Why counting and double negation do not remove the missing step}
\label{sec:logic}
\subsection{A precise counterexample to a general shortcut}
Let $R(e,t)$ say that deterministic machine $e$ on input $e$ first halts
at time $t$. Bounded simulation and a finite prefix check make $R$
primitive recursive. Each fiber has at most one member.

There is no total computable function $b(e)$ bounding that member.
Otherwise simulate $e(e)$ for $b(e)$ steps. A halt answers yes.
If the machine ever halted later, its first time would violate the bound;
thus failure to halt in that prefix answers no. This would decide diagonal
halting. Its impossibility follows directly: given a purported decider
$D$, construct $M(e)$ to loop when $D(e)$ says yes and halt when it
says no; evaluating it at its own machine code contradicts either answer.

So a decidable predicate may have the uniform cardinality bound one
without a uniform computable complete list or numerical bound.
These predicates are not asserted to be rational-point sets of curves.
The example refutes a general conversion from counting to enumeration,
not the possibility of an effective Mordell theorem.

At a simpler finite level, the empty predicate and $\{B+1\}$ agree
through any fixed cutoff $B$, both have cardinality at most one, and
differ on existence. This exact statement is checked in \texttt{Logic.lean}.

\subsection{A double-negated bound can exist without a positive witness}
For any predicate $P$ with at most one member, put
$\mathcal B(P)=\exists B\,\forall n(P(n)\to n\leq B)$.
Then intuitionistically
\[
                      \neg\neg\mathcal B(P).             \tag{12}
\]
Assume $\neg\mathcal B(P)$. A member $k$ would itself bound every member
by uniqueness, a contradiction. Hence $P$ is empty. But $0$ then bounds
it, again contradicting the assumption. This proves (12) without choosing
either a member or an emptiness proof.

Classical PRA may split on the unbounded statement $\exists n\,P(n)$
and prove the positive bound. For the halting fibers this gives a
classical uniform $\forall\exists\forall$ theorem with no computable
Skolem witness. Thus classical $\Pi^0_2$ conservation cannot be applied
directly to the Faltings-shaped formula. Intuitionistic arbitrary-matrix
witnessing is the essential ingredient above.

Negative translation keeps translated hypotheses and conclusions.
It does not both lower an arbitrary strong proof to $\iPRA$ and
remove double negations around the bound. Similarly, a finite prefix
search cannot certify the universal completeness statement
$\HB(\mathbf a,B)$ merely by checking points below $B$.

A proof-certificate checker is a valid alternative interface only after
its soundness and the existence of accepted certificates have been
proved. In the general case, that existence requirement retains the
unsolved finiteness work. It cannot be used as a hidden assumption.

\section{An unconditional genus-three instance: the Fermat quartic}
\label{sec:fermat-instance}

The general finiteness theorem remains unresolved in the requested setting.
There is nevertheless a substantial fixed-curve instance for which the
entire arithmetic argument can be given by primitive recursive operations
and quantifier-free induction.  We prove
\[
 C:\ X^4+Y^4=Z^4,
 \qquad
 C(\mathbb Q)=
 \{[1:0:1],[-1:0:1],[0:1:1],[0:-1:1]\}.
 \tag{FQ}\label{eq:fq-points}
\]
This is a classification, with a height bound and an explicit complete list.
It is stronger than merely ruling out arbitrarily long lists of distinct
points.  The arithmetic proof does not invoke Faltings' theorem.

The elementary descent below follows the mathematical route also recorded
in the inspected \texttt{diophantine\_geometry\_ipra.zip} draft, but every
ingredient needed here is proved.  The argument is a human-readable
arithmetic proof; this section does not claim that its deeply embedded
\texttt{Prf PRAAx} proof object has already been constructed.  The ordinary
geometric identification of this particular curve as a smooth genus-three
curve is separated at the end.

\subsection{Finite arithmetic conventions}

Signed integers are represented by a sign and a magnitude, with a unique
code for zero.  A rational number is represented by a reduced pair
$(u,v)$ with $u\in\mathbb Z$, $v\in\mathbb N$, $v>0$ and
$\gcd(|u|,v)=1$.  A rational projective point is a nonzero rational triple
modulo multiplication by a nonzero rational.  Its canonical code is an
integer triple $(X,Y,Z)$ with
\[
 \gcd(|X|,|Y|,|Z|)=1
 \quad\hbox{and first nonzero coordinate positive}.
 \tag{F1}
\]
Here the three-variable gcd means the gcd of the first two magnitudes
followed by the third.  Multiply the rational coordinates by the product
of their positive denominators, divide the resulting integers by their
positive gcd, and fix the sign.  This produces (F1) by a specified finite
algorithm.  Homogeneity preserves the quartic equation through these
operations.

For completeness, this representative is unique.  Suppose two primitive
integer triples differ by a rational factor $u/v$, with $v>0$ and
$\gcd(|u|,v)=1$.  From $vP_i=uQ_i$, Euclid's divisibility lemma, proved
below, implies that $v$ divides every $Q_i$.  Primitivity gives $v=1$.
The same argument in the other direction gives $|u|=1$, and the sign
convention gives $u=1$.  Every assertion here reduces to integer arithmetic
on three entries.  Define the nonlogarithmic height of a canonical point by
\[
 H(X,Y,Z)=\max(|X|,|Y|,|Z|).
\]

The calculations use addition, multiplication, comparison, parity, bounded
division, bounded gcd, and bounded square-root searches.  All are totalized
on malformed inputs, for example by returning zero.  A division $a/b$ in a
valid branch below is used only after $b>0$ and $b\mid a$ have been proved.
It can be implemented by testing $q\leq a$ for $bq=a$.  An exact square
root bounded by $N$ is obtained by testing the integers $r\leq N$ for
$r^2=a$.  Thus no unbounded minimization is concealed in the notation.

\subsection{The gcd and square-splitting lemmas}

The zero branches needed for projective normalization are fixed explicitly:
$\gcd(a,0)=a$, $\gcd(0,b)=b$, and $\gcd(0,0)=0$ for nonnegative
inputs. When one input is positive and the other is zero, the Bezout
coefficient pairs are respectively $(1,0)$ and $(0,1)$, and the
common-divisor properties follow immediately. A nonzero triple therefore
has a positive three-variable gcd even when two coordinates are zero.
The remaining computation is the positive-input case proved next.

\begin{lemma}[Bounded gcd and its certificates]\label{lem:fq-gcd}
For positive integers $a,b$ one can compute $g>0$ and signed integers
$u,v$ such that
\[
 g\mid a,\qquad g\mid b,\qquad ua+vb=g,
\]
and every common divisor of $a,b$ divides $g$.  These assertions have
primitive recursive witnesses and proofs by quantifier-free induction.
\end{lemma}

\begin{proof}
Start with $(a_0,b_0)=(\max(a,b),\min(a,b))$.  If $b_i>0$, perform bounded
division $a_i=q_i b_i+r_i$, $0\leq r_i<b_i$, and set
$(a_{i+1},b_{i+1})=(b_i,r_i)$.  If $b_i=0$, leave the pair fixed.
Before it stops, the positive sum $a_i+b_i$ strictly decreases: indeed
$b_i+r_i<a_i+b_i$.  Hence $a+b+1$ iterations suffice to reach a zero
second entry.  The first entry is positive throughout the stopped state.

Track both entries as signed linear combinations of the original inputs.
At a division step, replace the coefficient pair for the first entry by
that for the second, and replace the coefficient pair for the second by
the old first pair minus $q_i$ times the old second pair.  These rules
preserve the two displayed linear-combination equations.  If
$M=\max(a,b)$, then $q_i\leq M$, and every tracked coefficient after $j$
steps has magnitude at most $(M+1)^{j+1}$.  In particular, a finite storage
bound is available, although primitive recursion does not need one for
its definition.

For any fixed divisor $d$, the equations
$a_i=q_i b_i+r_i$ and $r_i=a_i-q_i b_i$ show that
\[
 (d\mid a_i\ \wedge\ d\mid b_i)
 \ \longleftrightarrow\
 (d\mid b_i\ \wedge\ d\mid r_i).
\]
At the stopped pair $(g,0)$ this proves both that every original common
divisor divides $g$ and, on taking $d=g$, that $g$ divides both original
inputs.  The coefficient invariant supplies $u,v$.

All induction invariants used in this argument are quantifier-free once
divisibility is replaced by its bounded characteristic function.  The
divisor $d$ is a free numerical parameter during induction, not an
unbounded quantifier in its induction formula.  Termination can use the
quantifier-free invariant
$b_i=0\ \vee\ (a_i+b_i+i\leq a+b)$ and the fixed fuel $a+b+1$.
\end{proof}

If $g=\gcd(a,b)$, the positive integers $a/g$ and $b/g$ are coprime:
a common divisor $d$ would make $gd$ a common divisor of $a,b$, hence
$gd\mid g$ and $d=1$.  This supplies the gcd normalization used below.

\begin{lemma}[Coprime cancellation and square extraction]
\label{lem:fq-squares}
The following implications have primitive recursive witness operations.
\begin{enumerate}
\item If $\gcd(a,b)=1$ and $a\mid bc$, then $a\mid c$.
\item If $d>0$ and $d^2\mid c^2$, then $d\mid c$.
\item If $a,b,c>0$, $\gcd(a,b)=1$, and $ab=c^2$, then there are
positive integers $r,s$ with $a=r^2$, $b=s^2$, and $c=rs$.
\end{enumerate}
\end{lemma}

\begin{proof}
For (1), multiply a Bezout identity $ua+vb=1$ by $c$.
Both terms on the right of $c=uac+vbc$ are divisible by $a$.
We also need two consequences of Bezout identities.  Squaring
$ua+vb=1$ and collecting terms gives
\[
 1=a(u^2a+2uvb)+v^2b^2,
\]
so coprimality of $a,b$ implies coprimality of $a,b^2$.
Multiplying a Bezout identity for $a,b$ by one for $a,c$ shows that
coprimality of $a$ with each of $b,c$ implies coprimality of $a,bc$.

For (2), the case $c=0$ is immediate.  Otherwise set
$h=\gcd(c,d)$, $c=hr$, and $d=hs$.  The normalized pair $r,s$ is
coprime.  Cancelling $h^2$ in the divisibility assumption gives
$s^2\mid r^2$, hence $s\mid r^2$.  But $s$ is coprime to $r^2$ by
the preceding Bezout calculation.  Applying (1), or its same Bezout
identity, gives $s\mid1$, so $s=1$.  Therefore $d=h\mid c$.

For (3), put $h=\gcd(a,c)$, $a=hu$ and $c=hv$.
Then $u,v$ are positive and coprime, and cancellation in $ab=c^2$
gives $ub=hv^2$.  Since $u$ is coprime to $v^2$, (1) implies
$u\mid h$.  Write $h=ue$.  Substitution and cancellation now yield
\[
 a=eu^2,\qquad b=ev^2,\qquad c=euv.
\]
The positive integer $e$ divides both $a,b$, so $e=1$.  Taking
$r=u$ and $s=v$ proves the result.  The construction uses only gcd
and exact division.  In particular, it is a witness-producing proof,
without an appeal to the existence of prime factorizations.
\end{proof}

The last lemma can be applied repeatedly to a product of any fixed number
of pairwise coprime positive factors.  Each factor is coprime to the
product of the others by the Bezout product calculation, so a square
product of four such factors makes all four factors squares.  This is
three uses of the lemma, not an unbounded search for square factors.

\subsection{Normalization and primitive Pythagorean parameters}

Call $(a,b,c,s)$ a square-area triangle if all four entries are positive
integers and
\[
 a^2+b^2=c^2,\qquad ab=2s^2.
 \tag{F2}\label{eq:fq-area}
\]
It is primitive if $\gcd(a,b)=1$.

\begin{lemma}[Primitive normalization]\label{lem:fq-normalize}
Every square-area triangle gives a primitive square-area triangle
$(A,B,C,S)$ with $C\leq c$ by a primitive recursive construction.
\end{lemma}

\begin{proof}
Put $g=\gcd(a,b)$, $A=a/g$, and $B=b/g$.  Then
$g^2\mid c^2$ by the Pythagorean equation, so
Lemma~\ref{lem:fq-squares}(2) gives $g\mid c$.  Set $C=c/g$.
We have $A^2+B^2=C^2$ and $\gcd(A,B)=1$.
The legs cannot both be even.  They cannot both be odd either, since
squares modulo $4$ are $0$ or $1$, whereas the sum of two odd squares
is $2$ modulo $4$.  Exactly one of $A,B$ is therefore even.
Consequently $K=AB/2$ is an integer, and the second equation in
\eqref{eq:fq-area} becomes $s^2=g^2K$.  The same square-divisibility
lemma gives $g\mid s$.  Set $S=s/g$ and cancel $g^2$ in both equations.
All output entries are positive, and $C=c/g\leq c$.
\end{proof}

\begin{lemma}[Explicit Pythagorean parameters]\label{lem:fq-pyth}
If $A,B,C>0$, $A^2+B^2=C^2$, and $\gcd(A,B)=1$, exchange the legs
if necessary to make $A$ even.  There are positive integers $m>n$ with
\[
 A=2mn,\qquad B=m^2-n^2,\qquad C=m^2+n^2,
 \tag{F3}\label{eq:fq-pyth}
\]
such that $\gcd(m,n)=1$ and $m,n$ have opposite parity.
The parameters are computable by searches bounded by $C$.
\end{lemma}

\begin{proof}
The parity argument just given makes $B$ and $C$ odd.
Also $B,C$ are coprime.  For if $d$ divides them both, then
$d^2\mid C^2-B^2=A^2$, so $d\mid A$ by
Lemma~\ref{lem:fq-squares}(2), and $d=1$.
Since $A>0$, the equation gives $C>B>0$.  Thus
\[
 U=(C+B)/2,\qquad V=(C-B)/2
\]
are positive integers with $U>V$.  A common divisor of $U,V$
divides $U+V=C$ and $U-V=B$, hence equals $1$.  Moreover
\[
 UV=(C^2-B^2)/4=(A/2)^2.
\]
Lemma~\ref{lem:fq-squares}(3) gives $U=m^2$, $V=n^2$, and
$A/2=mn$.  These equations give \eqref{eq:fq-pyth} and $m>n$.
A common divisor of $m,n$ would divide $U,V$, so they are coprime.
Their squares sum to the odd number $C$, so their parities differ.
Because $U,V\leq C$, their nonnegative square roots are at most $C$.
They are unique: if $r<t$ are nonnegative, then
$t^2-r^2=(t-r)(t+r)>0$.  A scan of $0,\ldots,C$ therefore returns
the required roots and has a proved success condition.
\end{proof}

\subsection{A descending map, with a numerical stopping bound}

\begin{lemma}[One explicit descent step]\label{lem:fq-step}
A primitive square-area triangle $(A,B,C,S)$ gives another primitive
square-area triangle $(A',B',C',S')$ with $0<C'<C$.
The step is a total primitive recursive function after assigning zero
output to invalid inputs.
\end{lemma}

\begin{proof}
Use the ordering and parameters of Lemma~\ref{lem:fq-pyth}.
Substituting \eqref{eq:fq-pyth} in $AB=2S^2$ and cancelling $2$
gives
\[
 mn(m+n)(m-n)=S^2.
 \tag{F4}
\]
All four factors are positive and pairwise coprime.
For pairs involving $m$ or $n$, subtraction reduces a common divisor
to a common divisor of $m,n$.  For the remaining pair, a common divisor
of $m+n,m-n$ divides $2m,2n$, hence divides $2$ by a Bezout identity
for $m,n$.  Both $m+n,m-n$ are odd, so that divisor is $1$.
Square splitting supplies positive integers $r,s,t,u$ such that
\[
 m=r^2,\qquad n=s^2,\qquad m+n=t^2,\qquad m-n=u^2.
 \tag{F5}\label{eq:fq-four-squares}
\]
The integers $t,u$ are odd, $t>u>0$, and they are coprime, since
any common divisor would divide the coprime pair $m+n,m-n$.
The integer $s$ is even.  Indeed, if $m$ were even and $n$ odd,
then $m=r^2\equiv0\pmod4$ and $n=s^2\equiv1\pmod4$,
making $m-n=u^2\equiv3\pmod4$, impossible.  Hence $n$ is even,
and so its square root $s$ is even.

Define
\[
 A'=(t+u)/2,\qquad B'=(t-u)/2,\qquad
 C'=r,\qquad S'=s/2.
 \tag{F6}\label{eq:fq-step}
\]
These are positive integers.  Direct expansion gives
\[
 (A')^2+(B')^2=(t^2+u^2)/2=m=r^2=(C')^2,
\]
and
\[
 A'B'=(t^2-u^2)/4=n/2=s^2/2=2(S')^2.
\]
A common divisor of $A',B'$ divides their sum $t$ and their
difference $u$, so the new triangle is primitive.  Finally,
\[
 0<C'=r\leq r^2=m<m^2+n^2=C.
 \tag{F7}\label{eq:fq-decrease}
\]
The last inequality follows from $m\geq1$, $n\geq1$ and
$m^2\geq m$.

For a completely specified algorithm, first test primitivity and both
equations.  On failure return $(0,0,0,0)$.  On success, exchange the
legs to make the first even, find $m,n$ as the unique exact roots of
$(C+B)/2,(C-B)/2$, and find $r,s,t,u$ as the exact roots in
\eqref{eq:fq-four-squares}.  Return \eqref{eq:fq-step}.
Every root search can use the input bound $C$: both $m,n\leq C$,
and $m+n\leq m^2+n^2=C$; the other radicands are no larger.
All denominators in the valid branch are the numeral $2$ and have
just been proved to divide their numerators.  Exact-root failure can
also be totalized by returning zero, and the proof shows that this
failure branch is never taken on a valid input.
\end{proof}

\begin{theorem}[No square-area triangle in intuitionistic PRA]
\label{thm:fq-no-area}
There are no positive integers satisfying \eqref{eq:fq-area}.
\end{theorem}

\begin{proof}
By Lemma~\ref{lem:fq-normalize}, a supposed solution produces a
primitive one, encoded by a state $q$.  Let $G(q)$ be the Boolean
predicate of being such a primitive solution, let $D$ be the total
step function from Lemma~\ref{lem:fq-step}, and let $\rho(q)$ be its
hypotenuse coordinate.  Define
\[
 q_0=q,\qquad q_{i+1}=D(q_i)
\]
by primitive recursion.  Under the assumption $G(q)$, induction gives
\[
 G(q_i)\quad\wedge\quad \rho(q_i)+i\leq\rho(q).
 \tag{F8}\label{eq:fq-invariant}
\]
The base is immediate.  For the step, Lemma~\ref{lem:fq-step} gives
$G(q_{i+1})$ and
$\rho(q_{i+1})+1\leq\rho(q_i)$; adding $i$ proves the inequality
at $i+1$.  This is quantifier-free induction: $G$ and the comparison
are fixed primitive recursive Boolean tests, and the state iteration
is a fixed primitive recursive function.

Now instantiate \eqref{eq:fq-invariant} at the specified numeral-valued
term $i=\rho(q)$.  Its inequality forces $\rho(q_i)=0$, whereas
$G(q_i)$ requires $\rho(q_i)>0$.  The contradiction discharges the
assumption $G(q)$.  The argument uses exactly the initial hypotenuse
as descent fuel and does not select a least counterexample or an
infinite sequence.  Negation is proved by contradiction from the
assumption being negated; no double-negation elimination is used.
\end{proof}

\subsection{From descent to the complete quartic list}

\begin{lemma}[Fourth-power difference]\label{lem:fq-difference}
There are no positive integers $x>y$ and $z$ with $x^4-y^4=z^2$.
Consequently, $a^4+b^4=c^4$ has no solution in positive integers.
\end{lemma}

\begin{proof}
A solution of $x^4-y^4=z^2$ would produce the positive integer
triangle
\[
 a=2x^2y^2,\qquad b=x^4-y^4,\qquad c=x^4+y^4,
 \qquad s=xyz.
\]
Expanding shows
\[
 (2x^2y^2)^2+(x^4-y^4)^2=(x^4+y^4)^2,
 \qquad ab=2x^2y^2z^2=2s^2,
\]
contradicting Theorem~\ref{thm:fq-no-area}.
If $a,b,c>0$ satisfy $a^4+b^4=c^4$, then $c>a$ and
$c^4-a^4=(b^2)^2$, a forbidden fourth-power difference.
\end{proof}

\begin{theorem}[The exact rational points of the Fermat quartic]
\label{thm:fq-complete}
In canonical integer projective codes, intuitionistic PRA proves that
every point of $X^4+Y^4=Z^4$ is one of
\[
 (1,0,1),\quad (1,0,-1),\quad (0,1,1),\quad (0,1,-1).
 \tag{F9}\label{eq:fq-canonical}
\]
Each listed triple is a point, the triples are pairwise distinct, and
every point has height $1$.  Therefore \eqref{eq:fq-points} holds.
\end{theorem}

\begin{proof}
Take a nonzero canonical integer triple satisfying the equation.
If $Z=0$, nonnegativity of fourth powers gives $X=Y=0$, a
contradiction.  Thus $|Z|>0$.  Equality of integers is decidable.
If both $X$ and $Y$ were nonzero, the positive integers
$|X|,|Y|,|Z|$ would contradict Lemma~\ref{lem:fq-difference}.
Case distinction on $X=0$ and $Y=0$ therefore leaves a zero
coordinate among $X,Y$.

If $Y=0$, the equation is $|X|^4=|Z|^4$.  Strict monotonicity of
the fourth power on nonnegative integers gives $|X|=|Z|$.
Primitivity forces this common positive value to be $1$.  The
first-nonzero-positive convention gives $X=1$, and $Z=1$ or $Z=-1$.
These are the first two triples of \eqref{eq:fq-canonical}.
If $X=0$, the identical argument gives $Y=1$ and $Z=1$ or $Z=-1$,
the last two triples.  They exhaust all cases.

Substitution verifies the equation for each listed triple, and direct
coordinate comparison proves their distinctness and their height $1$.
The second canonical triple represents $[-1:0:1]$, and the fourth
represents $[0:-1:1]$, by multiplication by $-1$.
All uses of case distinction concern explicit integer comparisons;
every induction above has a quantifier-free Boolean invariant.
\end{proof}

\subsection{The usual genus-three interpretation}

This paragraph identifies the arithmetic equation with a curve in the
ordinary statement of Faltings' theorem.  It is not an extra premise in
the integer classification just proved.

Over any characteristic-zero field, the three partial derivatives of
$F=X^4+Y^4-Z^4$ are $4X^3,4Y^3,-4Z^3$.  They cannot vanish
simultaneously at a projective point: their vanishing would force
$X=Y=Z=0$.  The projective Jacobian criterion therefore makes the
quartic smooth.  Absolute irreducibility can also be checked directly.
Over any extension field $E$, dehomogenize to
$f(X,Y)=Y^4+X^4-1$.  If $f=A B$ were a nontrivial factorization in
$E[X,Y]$, the leading coefficients in $Y$ would multiply to $1$.
After scaling, both factors would be monic in $Y$ and of positive
$Y$-degree; a factor of $Y$-degree zero would be a unit.
Reducing at $X=1$ gives $A(1,Y)B(1,Y)=Y^4$.  Both monic factors
of positive degree must have zero constant term, since the only
monic divisors of $Y^4$ are powers of $Y$.  Hence $(X-1)$ divides
both $A(X,0)$ and $B(X,0)$, so $(X-1)^2$ divides $X^4-1$.
This is impossible because
\[
 X^4-1=(X-1)(X^3+X^2+X+1),
 \qquad 1+1+1+1=4\ne0.
\]
A homogeneous factor which becomes a unit after setting $Z=1$
would be a power of $Z$ times a scalar, but $Z$ does not divide $F$.
Thus the homogeneous quartic is geometrically integral as well.

Its genus is $3$.  One can see the numerical calculation from the
Hilbert polynomial: for $n\geq4$, multiplication by $F$ is injective
on degree $n-4$ homogeneous polynomials, so the degree $n$ part of
$E[X,Y,Z]/(F)$ has dimension
\[
 \binom{n+2}{2}-\binom{n-2}{2}=4n-2=4n+1-3.
\]
For a smooth geometrically integral projective plane curve, the genus
is the number $g$ in this polynomial $dn+1-g$.  Equivalently, the
plane-curve genus formula gives $(4-1)(4-2)/2=3$; the usual geometric
theorem and its cohomological proof are recorded in the Stacks Project,
Section~53.9, Lemma~53.9.3, at
\url{https://stacks.math.columbia.edu/tag/0BYA}.
Thus \eqref{eq:fq-points} is a genuine genus-greater-than-one instance,
while its arithmetic proof requires none of this general geometry.

\subsection{Concrete Lean completion targets for this instance}

The proof gives a finite development with no missing finiteness theorem.
Implement the signed-number and tuple codes, and each fixed arithmetic
operation, as terms in the repository's \texttt{PR} syntax.  Prove
the bounded Euclidean and exact-root invariants, then the three
square lemmas, primitive normalization, the parameter equations,
and the preservation and decrease theorems for the displayed map $D$.
The required invariant for \texttt{PRA.qfInd} is exactly
\eqref{eq:fq-invariant}.  The final classification is a universal
formula whose matrix is a Boolean validity implication followed by
the four explicit coordinate equalities in \eqref{eq:fq-canonical}.

The desired final objects are an actual derivation
\[
 []\vdash[\mathtt{PRAAx}]\ \varphi_{\rm quartic}
\]
and its image under \texttt{pra\_iff\_native}.  The repository's
ambient theorem \path{LeanProofs.fermat_four_no_positive_nat_solutions}
is useful for comparison, but its existence does not construct that
object-theory derivation.  Conversely, completing this specific
arithmetic development would not complete the general Faltings
formalization: the explicit descent map is special to this quartic.


\section{An empty hyperelliptic family}
Fix $g\geq2$ and consider the smooth completion of
\[
                       y^2=-(x^{2g+2}+1).               \tag{13}
\]
With $x=u/v$, $y=r/s$, $v,s>0$, clearing denominators and writing
$a=|u|$, $c=|r|$ gives
\[
 c^2v^{2g+2}+s^2a^{2g+2}+s^2v^{2g+2}=0.
\]
All summands are nonnegative and the last is positive, a contradiction.
Power positivity is proved by quantifier-free induction on the exponent.
The other chart, $t=1/x$, $z=y/x^{g+1}$, has
$z^2=-(1+t^{2g+2})$. At infinity it would require $z^2=-1$,
which is impossible by the same square argument.
The complete rational-point list is therefore empty.

For its ordinary geometric interpretation, $X^{2g+2}+1$ is squarefree:
a common root with its derivative $(2g+2)X^{2g+1}$ would be zero,
not a root of the original polynomial. The smooth double cover has
$2g+2$ simple branch points and genus $g$.
The arithmetic nonexistence proof needs no genus theorem.
Both arithmetic contradictions are checked in \texttt{Curves.lean}
with no kernel axiom dependencies.

\section{Repository audit and the actual proof interface}
The inspected snapshot is
\texttt{d72c8b24637e342816dc92d11bef337748683e51}.
The audit covered the ConciseIPRA README, selected core Lean modules,
two relevant Diophantine drafts, and number-theory source. It was not an
exhaustive content audit of the whole repository. Provenance is included.

\begin{center}
\begin{tabularx}{\textwidth}{@{}p{0.35\textwidth}X@{}}
\toprule
Item & Contribution and boundary\\ \midrule
\texttt{PR.lean} & PR syntax, defining equations, \texttt{PRAAx},
and quantifier-free induction.\\
\texttt{Deduction.lean} & Embedded first-order proof objects
\texttt{Prf Ax}, distinct from ambient Lean propositions.\\
\texttt{NativeBiInterpretation.lean} & Faithful translation of
already constructed derivations.\\
Diophantine-geometry draft & Explicit arithmetic and finite descent;
general Faltings is explicitly outside its claims.\\
\texttt{FermatFour.lean} & Ambient Mathlib-backed Fermat-four statements;
no automatic embedded $\iPRA$ derivation.\\
\bottomrule
\end{tabularx}
\end{center}

The exact bridge has the form
\[
 (\Gamma\vdash[\mathtt{PRAAx}]\varphi)
 \ \longleftrightarrow\
 \mathtt{NPrf}(\Gamma.\mathtt{map}\,\mathtt{NFm.ofPRA})
                    (\mathtt{NFm.ofPRA}\,\varphi).
\]
This is \texttt{pra\_iff\_native}. It requires a proof on its left.
An ambient Lean theorem about native \texttt{Nat} is not that proof.
Equiconsistency alone also transports no particular theorem.
No unrestricted reflection from Lean truth to $\mathtt{Prf\ PRAAx}$
is assumed.

\section{Lean definitions, checked results, and a completion plan}
\label{sec:lean}
\subsection{Delivered definitions}
The project uses Lean 4.32.0 and \texttt{Std}, without Mathlib.
\begin{description}
\item[\texttt{FiniteSearch.lean}]
Defines complete lists, code bounds, positive finiteness, cardinality
bounds, interval enumeration, and list maximum. It proves exact search
membership and positive finiteness iff boundedness.
\item[\texttt{Logic.lean}]
Defines at-most-one, boundedness, finite-prefix agreement, the empty
predicate, and a hidden singleton. It proves the double-negation and
finite-prefix separation lemmas.
\item[\texttt{Curves.lean}]
Defines quintic coefficients, homogeneous evaluation, canonical coordinate
tuples, points, component height, the Sylvester determinant, goodness,
height boundedness, and the unproved uniform goal. It implements
bounded enumeration and proves its exact specification and list/height
equivalence for the concrete point predicate. It also checks the sign
contradictions for (13).
\end{description}

The central concrete checked statement is
\[
 p\in\mathtt{boundedPoints}(f,B)
 \leftrightarrow
 \mathtt{AffinePoint}(f,p)\wedge\mathtt{height}(p)\leq B.
\]
The structure \texttt{UniformHeightData} requires an explicit bound
function and its correctness proof. Its implication to the goal is
conditional on receiving such data; no inhabitant is supplied.
An arbitrary Lean function is not automatically primitive recursive.

The determinant \emph{definition} is checked; a full semantic theorem
linking it to a library of smooth projective curves is not included.
The files do not develop schemes, differential sheaves, or
Riemann--Hurwitz.

\subsection{Verification and its precise limits}
All three modules compile. The audit logs report:
\begin{center}
\begin{tabularx}{\textwidth}{@{}p{0.56\textwidth}X@{}}
\toprule
Principal results & Kernel dependencies\\ \midrule
Generic enumeration and list/bound equivalence &
\texttt{propext}; empty-list equivalence has none.\\
Logical double-negation and prefix results & None.\\
Concrete quintic box and list/bound equivalence &
\texttt{propext}, \texttt{Quot.sound}.\\
Conditional uniform-data implication and sign contradictions & None.\\
\bottomrule
\end{tabularx}
\end{center}
There is no \texttt{Classical.choice}, \texttt{sorryAx}, or custom axiom
in those audited results. The extensionality and quotient dependencies
come through standard-library lemmas. This native kernel audit is not
a proof that the development is already embedded in $\iPRA$.

The Python program compares the cleared-denominator search with an
independent rational-evaluation and exact-square-root search on four
quintics at five bounds, including zero. All twenty comparisons pass.
It also checks small discriminants and malformed codes.
For $y^2=x^5-x+1$ at height at most $12$, it returns
\[
                  (-1,\pm1),\quad(0,\pm1),\quad(1,\pm1).
\]
The discriminant is $2869$. The output explicitly claims completeness
only within that box, and does not assert that those are all points.

\subsection{A complete route for the proved finitary arguments}
The following order gives a concrete object-theory implementation plan.
\begin{enumerate}
\item Implement the integer, tuple, and list encodings above in the
repository's PR constructors, with decoder equations and the cap (9).
\item Build the point, goodness, divisibility, and comparison
characteristic terms and prove their specifications.
\item Construct the \texttt{PRA.qfInd} proof of the Boolean invariant (7),
keeping the candidate point as a free parameter; generalize afterwards.
\item Construct the maximum-scan proof and perform intuitionistic
existential elimination and introduction to prove (8).
\item For the Fermat quartic, implement the bounded gcd/root routines
and the explicit descent map of Section \ref{sec:fermat-instance}.
Use its displayed quantifier-free fuel invariant and finish the four-case
classification.
\item Apply \texttt{pra\_iff\_native} after the object-theory derivation
exists. Verify the precise translated formula with every coding predicate.
\end{enumerate}
Every mathematical step for this fixed-curve plan is supplied in the
human argument. The remaining work is substantial formalization;
the archive does not contain the completed embedded quartic derivation.

\subsection{The unresolved general endpoint}
With the numerical characteristic terms fixed, the missing formula is
\[
 \Phi_5=\forall c\bigl(\mathtt{good}(c)=1\to
 \exists B\,\forall p\,[
 \mathtt{point}(c,p)=1\to
 \mathtt{le}(\mathtt{height}(p),B)=1]\bigr).
\]
A full positive solution must construct
$[]\vdash[\mathtt{PRAAx}]\Phi_5$. Producing the bound is the mathematical
gap; neither filtering nor translating the proof calculus supplies it.

The original theorem for arbitrary number fields needs further geometric
definitions: a field presentation by an irreducible polynomial and a
rational basis; homogeneous projective equations; smoothness and geometric
integrality; and a genus definition with its verified connection to the
presentation. Points may use the first nonzero projective coordinate as
$1$, with other coordinates represented in the rational basis of $K$.
Their reduced rational coefficient heights again have enumerable balls.

This describes an arithmetization strategy, not an established primitive
recursive checker for every geometric hypothesis. One may accept finite
geometric certificates instead, but must then prove how they are obtained
for every legitimate input. The monic-quintic family avoids that issue
and already contains the unresolved positive bound question.

\section{What remains mathematically open in this package}
Naming Shafarevich finiteness, Parshin's construction, or alternative
Diophantine-approximation and $p$-adic proof methods does not provide
their constructive weak-arithmetic implementation. A valid positive
proof must make its witnesses, termination bounds, and induction
principles explicit wherever they are used to obtain finiteness.

There are two separate next targets. The strong one is \eqref{eq:target}
with primitive recursive bound production and an arithmetic correctness
proof. The weak one is a uniform cardinality statement such as (C),
possibly based on the recent discriminant bounds. It has a different
constructive meaning. No $\iPRA$ derivation of the analytic and geometric
estimates in those papers has been supplied here; explicit constants
alone do not bound the proof-theoretic strength of their derivation.

The finite reductions and special instances are complete at the stated
human-proof level. The general theorem remains visibly unproved.
This is a limitation of the present result, not an assertion of
unprovability.

\appendix
\section{Reproducing the artifacts}
The main source is \texttt{faltings\_ipra.tex}, with the local input
\texttt{sections/fermat\_instance.tex}. On a standard TeX installation,
compile using the command
\begin{quote}\ttfamily latexmk -pdf faltings\_ipra.tex\end{quote}
The included
\texttt{tools/build\_pdf.py} also supports the installed source layout
used here.

Run \texttt{lake build} in \texttt{lean/} with Lean 4.32.0, or run Lean
directly on each module. Run \texttt{python3 python/verify.py} for the
twenty independent finite-search comparisons.
The example command \texttt{python3 python/bounded\_points.py --bound 12}
prints a bounded-search result with its scope.

The verification directory records compiler versions, successful
compilation, and kernel dependencies. Provenance records the inspected
revision and source paths. The SHA-256 manifest covers delivered files.
Toolchain binaries, external papers, and intermediate TeX products
are omitted.

\begin{thebibliography}{10}
\bibitem{Stoll}
Michael Stoll. \emph{Determining the Rational Points on a Curve of Genus 2
and Mordell--Weil Rank 1}. Proceedings of the International Congress of
Mathematicians 2026, Vol.~3, pp.~484--498.
\url{https://doi.org/10.1137/25M1804650}.

\bibitem{FM}
Fernando Ferreira and Ant\'onio Marques.
\emph{Extracting Algorithms from Intuitionistic Proofs}.
Mathematical Logic Quarterly 44 (1998), 143--160.
\url{https://doi.org/10.1002/malq.19980440202}.
Author manuscript:
\url{https://webpages.ciencias.ulisboa.pt/~fjferreira/extracting.pdf}.
The theorem attributed to Wehmeier on printed p.~14 permits arbitrary
matrices.

\bibitem{YYZ}
Jiawei Yu, Xinyi Yuan, and Shengxuan Zhou.
\emph{Quantitativity on the Number of Rational Points in the Mordell
Conjecture}. Preprint, version 1, February 2, 2026.
\url{https://arxiv.org/abs/2602.01820}.
In particular, Theorem 1.7 gives discriminant-dependent cardinality bounds.

\bibitem{GFP}
Natalia Garcia-Fritz and Hector Pasten.
\emph{Effective Mordell for curves with enough automorphisms}.
Preprint, version 2, April 27, 2025.
\url{https://arxiv.org/abs/2503.10443}.

\bibitem{AL}
Levent Alp\"oge and Brian Lawrence.
\emph{Conditional algorithmic Mordell}.
Preprint, August 21, 2024.
\url{https://arxiv.org/abs/2408.11653}.

\bibitem{Stacks}
The Stacks Project Authors.
\emph{Plane curves}, Section 53.9, especially Lemma 53.9.3.
\url{https://stacks.math.columbia.edu/tag/0BYA}.
Used for the ordinary geometric interpretation of the quartic's genus.

\bibitem{Repo}
{\raggedright
Vladimir Reshetnikov and contributors. \emph{ProveIt}.
Snapshot
\href{https://github.com/VladimirReshetnikov/ProveIt/commit/d72c8b24637e342816dc92d11bef337748683e51}
{\texttt{d72c8b24637e342816dc92d11bef337748683e51}}.
Relevant paths:\par
\path{Logic/PrimitiveRecursiveArithmetic/ConciseIPRA};\par
\path{docs/drafts/diophantine_geometry_ipra.zip};\par
\path{NumberTheory/DiophantineEquations/Lean/DiophantineEquations/FermatFour.lean}.
\par}
\end{thebibliography}
\end{document}

